The application of deep learning techniques to the automated classification of prostate cancer has been extensively investigated, with significant advances particularly in the use of whole-slide images (WSIs), patch-based strategies, and multi-scale aggregation mechanisms.

\cite{ref-bulten2020} developed a deep learning system for automated Gleason grading using a dataset of 5,759 biopsies from Radboud University Medical Center, stratified into training (4,712), validation (497), and test (535) sets. The method, based on an extended U-Net architecture with semi-supervised labeling, achieved a quadratic weighted kappa ($\kappa_{quad}$) of 0.918 on the internal test set. System validation also included an observer study, in which the model outperformed 10 out of 15 pathologists, as well as evaluation on an independent external dataset from the University Hospital of Zurich, achieving $\kappa_{quad} > 0.700$.

\cite{ref-nagpal2019} developed a two-stage deep learning system (InceptionV3 CNN and k-NN) for Gleason grading using data from TCGA and two medical centers in the United States. The model was trained on 112 million annotated patches extracted from 1,226 slides and validated on an independent cohort of 331 prostatectomy slides. The system achieved an accuracy of 0.70 against an expert reference standard, significantly outperforming the average accuracy of 0.61 obtained by a panel of 29 pathologists ($p = 0.002$), and demonstrated superior performance in clinical risk stratification.

Patch-based approaches have evolved rapidly.
\cite{ref-kott2019} conducted a pilot study to develop a deep learning algorithm specifically focused on prostate biopsies using ResNet-18. The model was trained and validated on 14,803 image patches extracted from 85 WSIs from 25 patients at a single institution. The system achieved 91.5\% accuracy for malignancy detection (AUC = 0.83) and 85.4\% accuracy for multiclass Gleason score classification, including benign, Gleason 3, Gleason 4, and Gleason 5.

\cite{ref-lucas2019} proposed an approach based on the Inception-v3 CNN architecture for the detection of Gleason patterns and the automated determination of Grade Groups in prostate biopsies. The study used a dataset comprising 96 histological tissue sections from 38 patients, each with detailed pixel-level annotations. The method generated probability maps to compute biopsy composition. The system achieved 92\% accuracy in malignancy detection and 90\% in distinguishing between low-grade ($GP \le 3$) and high-grade ($GP \ge 4$) patterns. In the final slide-level Grade Group classification, the model achieved a quadratic weighted kappa ($\kappa_{quad}$) of 0.70.

\cite{ref-hammouda2021} proposed a system for Gleason grading using a three-scale pyramidal CNN and a pixel-wise classification strategy with majority voting. The model was developed using 608 WSI images from Radboud University Medical Center, applying a balancing strategy to mitigate the stroma pattern's predominance. The approach, which integrates edge-enhancement preprocessing, achieved 94\% accuracy in the final determination of Grade Groups and demonstrated computational efficiency, operating with approximately 1 million parameters.

Other studies expanded the focus to full WSIs. \cite{ref-singhal2022} developed a deep learning system for biopsy grading based on a U-Net architecture with a ResNeXt50 encoder, using domain-agnostic training strategies and active learning to mitigate staining variations. The model was developed using 6,670 WSIs from three different institutions and demonstrated strong generalization performance, achieving 83.1\% accuracy and a quadratic weighted kappa ($\kappa_{\text{quad}}$) of 0.93 on an external dataset, while maintaining AUCs above 0.93 for clinical risk stratification across different centers.

Despite these advances, the field still faces limitations related to data scale and methodology. Many studies rely on small datasets~\cite{ref-kott2019,ref-lucas2019}, which compromises generalization. Even robust works focused on domain adaptation, such as \cite{ref-singhal2022}, place limited emphasis on architectural efficiency and on ordinal regression. Moreover, the potential of noise-filtering techniques and uncertainty-based curation remains underexplored compared to traditional classification and segmentation approaches.

In this context, the present work distinguishes itself by systematically evaluating multiple EfficientNet variants (B0, B3, and B7) for ISUP grading on the PANDA dataset, combining ordinal regression and focal loss to model the ordinal and imbalanced nature of the classes. We also explore optimized ensemble strategies that consistently outperform individual models and incorporate a Shannon entropy–based filtering step to mitigate noisy labels, an approach not explored in the studies reviewed in this work.
